%! Logarithmic Functions — Unit 2, Lesson 2.11 — Group Activity (Answer Key)
\documentclass[10pt]{article}
\usepackage{precalculus-article}
\usepackage{precalculus-key}   % pulls in -boxes; do NOT also load -boxes
\newcommand{\TallMath}[1]{$\displaystyle #1\rule[-1.4em]{0pt}{3.2em}$}

\begin{document}

\pageheader{Unit 2, Lesson 2.11}{Group Activity --- Answer Key}
\namepartnerperiod

\vspace{0.08in}

\begin{scenariobox}[Seismic Data Analysis]{sky}
A geologist uses $f(x) = \log_{10}(x)$ to interpret seismic intensity readings, where $x$ is
the measured wave intensity. Your group will analyze the characteristics of logarithmic functions
in this context and beyond.
\end{scenariobox}

\vspace{0.08in}

\begin{tcolorbox}[colback=white, colframe=black!40,
    title=\textbf{Tier R --- Remediate},
    fonttitle=\bfseries, arc=2mm, left=3mm, right=3mm, top=2mm, bottom=2mm]

\textbf{Work with $f(x) = \log_{10}(x)$.}

\begin{enumerate}[leftmargin=1.8em, itemsep=10pt]

\item State the domain and range of $f(x) = \log_{10}(x)$.

\vspace{4pt}
Domain: \ans{$x > 0$, i.e., $(0,\infty)$} \qquad Range: \ans{all reals, $(-\infty,\infty)$}

\vspace{4pt}
Equation of the vertical asymptote: $x = $ \ans{$0$}

\item Complete the table of values.

\vspace{4pt}
\begin{center}
\renewcommand{\arraystretch}{1.8}
\begin{tabular}{|c|c|}
\hline
\rowcolor{sky}
$x$ & $f(x) = \log_{10}(x)$ \\
\hline
$0.1$   & \ans{$-1$} \\
\hline
$1$     & \ans{$0$} \\
\hline
$10$    & \ans{$1$} \\
\hline
$100$   & \ans{$2$} \\
\hline
$1000$  & \ans{$3$} \\
\hline
\end{tabular}
\end{center}

\item The geologist records the readings $x = -5$, $x = 0$, and $x = 7.3$.
      Which readings are \textbf{invalid} inputs for $f$? Explain why.

\ansline{$x = -5$ and $x = 0$ are invalid. The domain of $\log_{10}(x)$ is $x > 0$;}
\ansline{you cannot take the log of a non-positive number.}

\item Based on your table, is $f$ increasing or decreasing on its domain?

\vspace{4pt}
$f$ is \ans{increasing} on $(0, \infty)$.

\begin{teachernote}
Students can observe from the table that as $x$ grows, $f(x)$ grows. This is consistent with
EK 2.11.A.2: logarithmic functions are always monotone and never have extrema.
\end{teachernote}

\end{enumerate}
\end{tcolorbox}

\vspace{0.08in}

\begin{tcolorbox}[colback=white, colframe=black!40,
    title=\textbf{Tier A --- Approaching Proficiency},
    fonttitle=\bfseries, arc=2mm, left=3mm, right=3mm, top=2mm, bottom=2mm]

\textbf{Work with $f(x) = 2\log_3(x)$.}

\begin{enumerate}[leftmargin=1.8em, itemsep=10pt]

\item State the domain and range. Write the equation of the vertical asymptote.

\vspace{4pt}
Domain: \ans{$(0,\infty)$} \quad Range: \ans{$(-\infty,\infty)$} \quad Asymptote: $x = $ \ans{$0$}

\item Fill in the end behavior using limit notation.

\vspace{4pt}
$\displaystyle\lim_{x \to 0^+} 2\log_3(x) = $ \ans{$-\infty$} \qquad
$\displaystyle\lim_{x \to \infty} 2\log_3(x) = $ \ans{$+\infty$}

\item Is $f$ increasing or decreasing? How do you know from the inverse relationship?

\ansline{$f$ is increasing. Its inverse $g(x) = 3^x$ ($b = 3 > 1$) is always increasing,}
\ansline{so the inverse logarithmic function is also always increasing (EK 2.11.A.2).}

\item Is the graph of $f$ concave up or concave down? Does $f$ have any local maximum or
      minimum on its natural domain? Explain using EK 2.11.A.2.

\ansline{Concave down. By EK 2.11.A.2, logarithmic functions have no local extrema on their}
\ansline{natural domain and are always concave in one direction without inflection points.}

\item Explain in words what happens to the output of $f(x) = 2\log_3(x)$ as $x$ gets
      very close to zero from the right.

\ansline{The output decreases without bound (goes to $-\infty$). The graph has a vertical}
\ansline{asymptote at $x = 0$, so $f(x) \to -\infty$ as $x \to 0^+$ (EK 2.11.A.4).}

\end{enumerate}
\end{tcolorbox}

\vspace{0.08in}

\begin{tcolorbox}[colback=white, colframe=black!40,
    title=\textbf{Tier E --- Extension},
    fonttitle=\bfseries, arc=2mm, left=3mm, right=3mm, top=2mm, bottom=2mm]

\textbf{Additive Transformation Identification Test.}

\begin{enumerate}[leftmargin=1.8em, itemsep=10pt]

\item The table below shows $f(x) = \log_3 x$. Check whether the inputs are proportional
      over equal-length output intervals by computing each ratio.

\vspace{4pt}
\begin{center}
\renewcommand{\arraystretch}{1.8}
\begin{tabular}{|c|c|c|}
\hline
\rowcolor{sky}
$x$ & $f(x)$ & Input ratio (consecutive) \\
\hline
$1$  & $0$ & --- \\
\hline
$3$  & $1$ & \ans{$3/1 = 3$} \\
\hline
$9$  & $2$ & \ans{$9/3 = 3$} \\
\hline
$27$ & $3$ & \ans{$27/9 = 3$} \\
\hline
$81$ & $4$ & \ans{$81/27 = 3$} \\
\hline
\end{tabular}
\end{center}

\vspace{4pt}
Are the inputs proportional? \ans{Yes} \quad Constant ratio: \ans{3}

\item Now form $g(x) = f(x + 2) = \log_3(x + 2)$. Find the new inputs by solving
      $\log_3(x + 2) = k$ for each output $k = 0, 1, 2, 3, 4$.

\vspace{4pt}
\begin{center}
\renewcommand{\arraystretch}{1.8}
\begin{tabular}{|c|c|c|}
\hline
\rowcolor{sky}
$x$ (new input for $g$) & $g(x)$ & Input ratio (consecutive) \\
\hline
\ans{$-1$} & $0$ & --- \\
\hline
\ans{$1$}  & $1$ & \ans{$1/(-1)= -1$} \\
\hline
\ans{$7$}  & $2$ & \ans{$7/1 = 7$} \\
\hline
\ans{$25$} & $3$ & \ans{$25/7 \approx 3.57$} \\
\hline
\ans{$79$} & $4$ & \ans{$79/25 \approx 3.16$} \\
\hline
\end{tabular}
\end{center}

\vspace{4pt}
Are the inputs of $g$ proportional? \ans{No} \quad This confirms EK 2.11.A.3.

\begin{teachernote}
For $g(x) = \log_3(x+2) = k$: $x + 2 = 3^k$, so $x = 3^k - 2$.
$k=0$: $x=-1$; $k=1$: $x=1$; $k=2$: $x=7$; $k=3$: $x=25$; $k=4$: $x=79$.
Ratios are not constant, confirming that the additive shift destroys the proportional-input property.
\end{teachernote}

\item A mystery function $h$ has the following table. Determine whether $h$ could be
      logarithmic by testing whether the inputs are proportional over equal output intervals.

\vspace{4pt}
\begin{center}
\renewcommand{\arraystretch}{1.8}
\begin{tabular}{|c|c|}
\hline
\rowcolor{sky}
$x$ & $h(x)$ \\
\hline
$2$  & $1$ \\
\hline
$8$  & $2$ \\
\hline
$32$ & $3$ \\
\hline
$128$& $4$ \\
\hline
\end{tabular}
\end{center}

\vspace{4pt}
Input ratios: \ans{$8/2 = 4$;\quad $32/8 = 4$;\quad $128/32 = 4$}

\vspace{4pt}
Is $h$ logarithmic? Justify. \ansline{Yes. The inputs are proportional (each multiplied by 4) over equal output intervals}
\ansline{(output increases by 1 each time), so $h$ is logarithmic --- specifically $h(x) = \log_4 x$.}

\end{enumerate}
\end{tcolorbox}

\end{document}
